\begin{abstract}
Memory is the most expensive and constrained resource in cloud
computing.  Operating systems provide transparent memory compression
(macOS Compressed Memory, Linux zswap), but these mechanisms operate
below the allocator and cannot exploit allocator-level structure such
as size classes, object liveness, or allocation-site patterns.

This paper presents \sys, a compression-aware memory allocator that
transparently compresses cold heap pages to reduce resident set size
(RSS).  \sys interposes on \texttt{malloc}/\texttt{free} via library
interposition, requiring no application changes.  Three key techniques
improve compression effectiveness: (1)~call-site arena segregation
routes allocations from the same call site to the same arena, placing
structurally similar objects on the same pages; (2)~zero-on-free
clears deallocated slots, replacing stale data with highly
compressible zero runs; and (3)~adaptive multi-algorithm compression
uses fast LZ4 for newly cold pages and escalates to zstd for
long-lived cold data.  A signal-based fault handler transparently
decompresses pages on access, with per-slot decompression contexts to
support concurrent faults.

Evaluation on five real-world applications shows that \sys reduces RSS
by 17--44\% with negligible throughput impact: memcached serves
requests at full throughput with 25\% lower RSS, SQLite reduces
in-memory database RSS by 39\%, and DuckDB analytics queries complete
with 30\% lower peak memory.  An ablation study across 17
configurations isolates each optimization's contribution and confirms
that call-site segregation, adaptive compression, and zero-on-free are
each independently beneficial.
\end{abstract}
